
\documentclass[9pt]{elife}

\usepackage{amsmath,amssymb,mathtools}
\usepackage{siunitx}
\usepackage{graphicx}
\usepackage{booktabs}
\usepackage{hyperref}
\usepackage{tcolorbox}
\usepackage{bm}

\title{MacroIR: Exact recursive likelihoods for time-averaged observations in continuous-time Markov models}

\author[1*]{Luciano Moffatt}
\affil[1]{Instituto de Qu\'imica F\'isica de los Materiales, Medio Ambiente y Energ\'ia (INQUIMAE), CONICET; Facultad de Ciencias Exactas y Naturales, Universidad de Buenos Aires, Ciudad de Buenos Aires, C1428EHA, Argentina}
\corr{lmoffatt@qi.fcen.uba.ar}

\begin{document}
\maketitle

\begin{abstract}
Time-averaged measurements are ubiquitous in biology, yet most likelihood methods assume instantaneous sampling, yielding biased inference when intrinsic kinetics outpace acquisition. We present \emph{MacroIR}, an exact, recursive likelihood for interval-averaged observations of continuous-time Markov models (CTMCs). By conditioning on both the start and end of each measurement window, MacroIR computes the interval mean and variance in closed form and performs Kalman-like updates in the original $k$-state space. Leveraging linear-operator identities, MacroIR collapses the apparent $k^2$ start--end expansion to standard $k\times k$ algebra. A score-based self-consistency test (mean score $\approx 0$, variance $\approx$ Fisher information) is satisfied by MacroIR but not by common midpoint or non-recursive approximations, establishing internal informational consistency. MacroIR resolves a longstanding problem in stochastic inference from averaged data and is broadly applicable beyond ion channels.
\end{abstract}

\section*{Introduction}
Many biological signals are averages over finite windows imposed by hardware integration or binning, whereas standard likelihoods for CTMCs implicitly assume instantaneous sampling. When intrinsic time constants are shorter than the acquisition window, midpoint or single-point approximations induce bias in kinetic parameters and model evidence. Here we introduce \emph{MacroIR}, a method that yields exact interval-averaged likelihoods with a recursive (Kalman-like) propagation of hidden-state moments. The approach is exact for true time integrals, and provides an internal validation via a score test whose empirical mean is $\approx 0$ and whose variance matches the Fisher information when the likelihood is correctly specified. We demonstrate the method on synthetic and experimental datasets and discuss the methodological scope, including the caveat that explicit analog filtering (e.g.\ Bessel) is not modeled unless the filter is included in the state-space, which increases computational cost.

A prior biological application used MacroIR to resolve gating asymmetry in P2X receptors (Moffatt, 2025; \emph{Communications Biology}, in press). The present work focuses on the method itself.

\section*{Theory}
\subsection*{Interval-averaged observation model}
Let $x_t$ be a CTMC on $k$ states with generator $Q$ and transition matrix $P(t)=e^{Qt}$. Let $\bm{\gamma}\in\mathbb{R}^k$ be a linear observable (e.g.\ state-wise current). For a window $[0,t]$ the measured signal is the average
\begin{equation}
\bar y_{0\to t} \;=\; \frac{1}{t}\int_0^t \gamma(x_s)\,ds.
\end{equation}
Conditioned on the start and end states $(i,j)$, the expected interval signal admits a spectral form
\begin{equation}
\label{eq:gammaij}
\gamma_{i\to j}(t) \;=\; \frac{1}{P_{i\to j}(t)} \sum_{n_1,n_2} V_{i n_1}\,(V^{-1}\Gamma V)_{n_1 n_2}\,V^{-1}_{n_2 j}\,E_2(\lambda_{n_1} t, \lambda_{n_2} t),
\end{equation}
where $Q=V\Lambda V^{-1}$, $\Gamma=\mathrm{diag}(\bm{\gamma})$, and $E_2(x,y)=(e^{x}-e^{y})/(x-y)$ (and $E_2(x,x)=e^x$). Equation~\eqref{eq:gammaij} implicitly marginalizes all internal paths without building a $k^2$ object.

\subsection*{Predicted mean and variance for one interval}
Let $(\bm{\mu}^{\mathrm{prior}}_0,\Sigma^{\mathrm{prior}}_0)$ be the mean and covariance of state occupancy at the interval start, and $N_{\mathrm{ch}}$ the number of independent replicas (e.g.\ channels). Define the interval-averaged first and second moment vectors $\bm{\gamma}_0$ and $\mathbf{E}$ assembled from $\gamma_{i\to j}(t)$ and $\gamma_{i\to j}^2(t)$. Then
\begin{align}
y^{\mathrm{pred}}_{0\to t} &= N_{\mathrm{ch}}\,(\bm{\mu}^{\mathrm{prior}}_0 \cdot \bm{\gamma}_0),\\
\sigma^2_{y^{\mathrm{pred}}_{0\to t}} &= \frac{\epsilon^2}{t} + \nu^2 + N_{\mathrm{ch}}\left[\bm{\gamma}_0^\top\!\big(\Sigma^{\mathrm{prior}}_0 - \mathrm{diag}\,\bm{\mu}^{\mathrm{prior}}_0\big)\bm{\gamma}_0 \;+\; \bm{\mu}^{\mathrm{prior}}_0 \cdot \mathbf{E}\right],
\end{align}
where $\epsilon^2$ and $\nu^2$ are white and pink noise contributions estimated from baseline.

\subsection*{Recursive update (Kalman-like)}
With $P(t)=e^{Qt}$ and an observed interval average $y^{\mathrm{obs}}_{0\to t}$,
\begin{align}
\bm{\mu}^{\mathrm{post}}_{t} &= \bm{\mu}^{\mathrm{prior}}_0 P(t) \;+\; \frac{y^{\mathrm{obs}}_{0\to t}-y^{\mathrm{pred}}_{0\to t}}{\sigma^2_{y^{\mathrm{pred}}_{0\to t}}}\,K^\mu_{0\to t},\\
\Sigma^{\mathrm{post}}_{t} &= \Sigma^{\mathrm{prior}}_{t} \;-\; \frac{1}{\sigma^2_{y^{\mathrm{pred}}_{0\to t}}}\,(K^\mu_{0\to t})^\top K^\mu_{0\to t},
\end{align}
where the gain is the exact cross-covariance divided by the interval variance,
\begin{equation}
K^\mu_{0\to t} \;=\; \mathrm{Cov}\!\big(x_t,\,\bar y_{0\to t}\big) \;\big/ \; \mathrm{Var}\!\big(\bar y_{0\to t}\big),
\end{equation}
and has a closed form depending only on $P(t)$, $\gamma_{i\to j}(t)$ and the prior moments. This update is mathematically equivalent to a Kalman gain over start--end pairs, but expressed directly in the $k$-dimensional space.

\begin{tcolorbox}[colback=gray!5,colframe=gray!50,title=Boundary-conditioned two-point states (``meta-states'')]
We define a \emph{boundary-conditioned two-point state} as the ordered pair $(i,j)$ describing that the chain started in $i$ and ended in $j$ over a duration $t$. The $k^2$ such pairs form a space whose generator is the Kronecker sum $Q\oplus Q$. Although this formal construction underlies the derivation, MacroIR never materializes it because $e^{(Q\oplus Q)t}=e^{Qt}\otimes e^{Qt}$ factorizes, so all boundary-conditioned interval statistics reduce to bilinear forms of $e^{Qt}$ and $\bm{\gamma}$ in the original $k$-dimensional space.
\end{tcolorbox}

\subsection*{Why no $k^2$ blow-up}
The joint start--end operator is $Q\oplus Q=Q\otimes I + I\otimes Q$, with exponential $e^{(Q\oplus Q)t}=e^{Qt}\otimes e^{Qt}$. Interval averages introduce time integrals of products of exponentials that yield kernels like $E_2$, which can be evaluated with standard $k\times k$ linear algebra and spectral identities. Because the observable is linear in state, all required expectations and covariances remain closed in the original space.

\subsection*{Scope}
MacroIR is exact for \emph{true} interval integrals. Incorporating explicit analog low-pass filtering (e.g.\ Bessel) is possible by augmenting the system with the filter state-space (Kalman/Bucy style) or by convolving the spectral kernels with the filter poles; both routes are substantially more expensive computationally.

% \begin{figure}[t]
% \centering
% \includegraphics[width=0.8\linewidth]{fig1.pdf}
% \caption{\textbf{Conceptual overview of MacroIR.} Interval averaging, boundary conditioning, and recursive update.}
% \label{fig:overview}
% \end{figure}

\section*{Methods}
\subsection*{Synthetic models and data generation}
We simulate $k$-state CTMCs with known $Q$ and observable $\bm{\gamma}$; time series are averaged over bins of duration $\Delta$ spanning sub- to supra-kinetic regimes. White/pink noise levels $(\epsilon^2,\nu^2)$ are matched to baseline segments. Priors are weakly informative on rates (log-uniform) and on noise parameters.

\subsection*{Inference settings and convergence}
Posterior exploration can use affine-invariant MCMC with parallel tempering. Convergence is monitored by potential scale reduction, effective sample sizes, and posterior concentration factors. Evidence is estimated via thermodynamic integration over the tempered ladder when required for model comparison.

\subsection*{Score test of sampling--likelihood duality}
Let $S(\theta)=\nabla_\theta \log \mathcal{L}(\theta;\text{data})$ be the score evaluated at the posterior mode or at known ground-truth $\theta^\star$ in simulations. Under correct likelihood specification,
\begin{equation}
\mathbb{E}[S(\theta^\star)] = 0, \qquad \mathrm{Var}(S(\theta^\star)) = \mathcal{I}(\theta^\star),
\end{equation}
where $\mathcal{I}$ is the Fisher information. We estimate the empirical mean and covariance of the score across intervals and compare them to a plug-in estimate of $\mathcal{I}$ computed from the MacroIR likelihood. In practice, MacroIR satisfies both properties, whereas midpoint and non-recursive approximations yield non-zero mean scores and information mismatches.

\section*{Results}
\subsection*{Synthetic recovery and bin-size dependence}
Across bin sizes, MacroIR recovers unbiased kinetic parameters and accurate uncertainty; midpoint/NR show bias and under-estimated variance at long bins. Posterior predictive checks confirm agreement of mean and variance of interval measurements under MacroIR.

% \begin{figure}[t]
% \centering
% \includegraphics[width=0.8\linewidth]{fig2.pdf}
% \caption{\textbf{Synthetic recovery vs bin size.} Bias (top) and variance (bottom) across methods.}
% \label{fig:synthetic}
% \end{figure}

\subsection*{Score-based self-consistency}
MacroIR yields empirical mean score $\approx 0$ and score variance matching the Fisher information within sampling error. Competing approximations fail these checks systematically, especially when $\Delta$ exceeds the fastest kinetic timescales.

% \begin{figure}[t]
% \centering
% \includegraphics[width=0.8\linewidth]{fig3.pdf}
% \caption{\textbf{Score test.} Mean score (left) and variance vs Fisher information (right) across methods.}
% \label{fig:score}
% \end{figure}

\subsection*{Experimental illustration}
On real macroscopic current recordings, MacroIR reproduces the amplitude and time course of responses and partitions variance into gating, white, and pink components. Evidence estimates favour the interval-exact recursive likelihood over non-recursive alternatives.

% \begin{figure}[t]
% \centering
% \includegraphics[width=0.8\linewidth]{fig4.pdf}
% \caption{\textbf{Experimental demonstration.} Posterior predictive fit and variance partitioning.}
% \label{fig:real}
% \end{figure}

\subsection*{Computational performance}
Runtime scales with $k$ and the number of intervals. In practice, the spectral evaluation of $E_2$ kernels and $k\times k$ multiplications keep costs manageable for the models considered; explicit analog filtering increases runtime by an order of magnitude.

% \begin{figure}[t]
% \centering
% \includegraphics[width=0.8\linewidth]{fig5.pdf}
% \caption{\textbf{Scaling.} Wall-clock vs $k$ and number of intervals.}
% \label{fig:scaling}
% \end{figure}

\section*{Discussion}
MacroIR provides an exact and efficient route to likelihoods for interval-averaged CTMC observations, resolving a key limitation of instantaneous-sample assumptions. The boundary-conditioned formulation ensures that temporal correlations induced by averaging are accounted for, and the score test gives an internal diagnostic of informational consistency. Although the present work treats true integrals, the framework can be extended to explicit analog filtering at additional computational cost. Beyond ion channels, MacroIR applies to any stochastic biological system observed through time-averaged windows (e.g.\ fluorescence trajectories, population dynamics, or electrophysiology).

\section*{Data and Code Availability}
All code implementing MacroIR, together with scripts to reproduce figures and simulations, will be made available in a public repository upon acceptance. Example notebooks and a reference C++/R implementation are planned.

\section*{Acknowledgements}
I thank colleagues for discussions that helped clarify the boundary-conditioned construction and the score-based validation. Computational resources provided by local HPC facilities are gratefully acknowledged.

\section*{Author Contributions}
LM designed the research, developed the theory, implemented the method, analyzed data, and wrote the manuscript.

\section*{References}
\begin{itemize}
  \item Moffatt L (2025) \emph{Bayesian inference of functional asymmetry in a ligand-gated ion channel}. Communications Biology, in press.
  \item Additional references will be added at submission (foundational CTMC texts, interval statistics, Kalman filtering, and related work).
\end{itemize}

\end{document}
